\section{1988: The Protocol, the Consecrations, and the First Rupture}

\subsection{The protocol of 5 May}

That an agreement was signed on 5 May 1988 by Cardinal Joseph Ratzinger,
Prefect of the Congregation for the Doctrine of the Faith, and Archbishop
Marcel Lefebvre is established by the Holy See's own acts, twice. The
motu proprio \work{Ecclesia Dei} of 2 July 1988 directs its new
commission to work for the full ecclesial communion of priests,
seminarians, communities or individuals linked to the Society ``in the
light of the Protocol signed on 5 May last by Cardinal Ratzinger and
Mons.\ Lefebvre''; and \work{Ecclesiae unitatem} of 2 July 2009 quotes
that clause verbatim twenty-one years later.\endnote{John Paul II, motu
proprio \work{Ecclesia Dei}, 2 July 1988, n.~6~a); Latin in \work{Acta
Apostolicae Sedis} 80 (1988) 1495--1498, at 1498 (``iuxta Protocollum
superiore die 5 mensis Maii obsignatum a Cardinali Ratzinger et ab
Archiepiscopo Lefebvre''), English on the Holy See website. Benedict
XVI, motu proprio \work{Ecclesiae unitatem}, 2 July 2009, n.~2, quoting
the same clause: \work{Acta Apostolicae Sedis} 101 (2009) 710--711, at
710. Both read 2026-07-25.}

The protocol's text is another matter. It was not printed in \work{Acta
Apostolicae Sedis} and was not located in any official Holy See
publication for this account. What its clauses contained---a doctrinal
declaration, a commission with Roman and Society members, provision for
one bishop---is reported in the literature and by both parties, but this
account does not assert the terms of a document it has not read in an
identified witness.\endnote{Bounded negative: a literal search of the
text layer of the scanned \work{Acta Apostolicae Sedis} 80 (1988) finds
the founder's name only in the papal letter of 8 April and in
\work{Ecclesia Dei}; the protocol is nowhere printed there, and no
official publication of its text was located on vatican.va (checked
2026-07-25). What is officially established is its date, its two
signatories, and the fact that the Holy See continued to treat it as the
reference point for reconciliation.}

What is certain is that the agreement did not hold. Within weeks
Archbishop Lefebvre resolved to consecrate bishops without a papal
mandate, and the Holy See moved from negotiation to warning.

\subsection{Warning, consecration, declaration}

On 16 June 1988 the Holy See published an information note in
\work{L'Osservatore Romano}; \work{Ecclesia Dei} cites it in its first
footnote as the record of the efforts made in the preceding
months.\endnote{\work{Ecclesia Dei}, footnote 1: ``Cf.\ Nota informativa
diei 16 Iunii 1988 in diurnario \guillemotleft L'Osservatore
Romano\guillemotright{} (17 Iunii 1988), pp.~1--2'' (\work{Acta
Apostolicae Sedis} 80 (1988) 1495 n.~1). Bounded negative: the text of
that note was not located in any digital witness reached for this
account. Its existence, date, and place of publication are officially
attested by the footnote itself.} On 17 June 1988 the Cardinal Prefect of
the Congregation for Bishops sent Lefebvre and the four priests a formal
canonical warning---the \latin{monitio} that penal law requires before a
declared penalty. That warning, its date, and its author are attested in
the body of \work{Ecclesia Dei}.\endnote{\work{Ecclesia Dei}, n.~3
(English, Holy See website): ``notwithstanding the formal canonical
warning sent to them by the Cardinal Prefect of the Congregation for
Bishops on 17 June last.''}

On 30 June 1988, at Écône, Archbishop Lefebvre consecrated four bishops:
Bernard Fellay, Bernard Tissier de Mallerais, Richard Williamson, and
Alfonso de Galarreta. \work{Ecclesia Dei} names all four in its third
paragraph. The co-consecrator was Antônio de Castro Mayer, bishop
emeritus of Campos in Brazil; his participation is attested officially not
by \work{Ecclesia Dei}, which does not name him, but by the 1996 note of
the Pontifical Council for Legislative Texts, which describes the
excommunication as concerning ``Bishop Marcel Lefebvre, the four bishops
ordained by him, and Bishop emeritus Antônio de Castro.''\endnote{Pontifical
Council for Legislative Texts, explanatory note \work{Sulla scomunica per
scisma in cui incorrono gli aderenti al movimento del Vescovo Marcel
Lefebvre}, 24 August 1996, covering letter (Italian text on the Holy See
website; published in \work{Communicationes} 29 [1997] 239--243), read
2026-07-25.}

The Holy See responded in two acts on successive days, and the difference
between them is important.

\begin{itemize}
\item On \textbf{1 July 1988} the Prefect of the Congregation for Bishops
issued a decree---known from its opening words as \work{Dominus Marcellus
Lefebvre}---formally declaring the excommunication \latin{latae
sententiae} already incurred. Its existence, date, author, and declaratory
character are attested by three later Roman documents; its text was never
printed in \work{Acta Apostolicae Sedis} and was not located in any
official witness here.\endnote{Attestations: Pontifical Council for
Legislative Texts, note of 24 August 1996, n.~1 (``il Decreto
\guillemotleft Dominus Marcellus Lefebvre\guillemotright{} della
Congregazione per i Vescovi, del 1 luglio 1988''); Secretariat of State,
note \work{Notandum circa excommunicationem quattuor Praesulum
Fraternitatis Sancti Pii X}, 4 February 2009, n.~1, \work{Acta
Apostolicae Sedis} 101 (2009) 145--146 (the penalty incurred on 30 June
1988 and ``formally declared on 1 July of the same year''); Congregation
for Bishops, decree of 21 January 2009, \work{Acta Apostolicae Sedis} 101
(2009) 150--151 (``la scomunica latae sententiae formalmente dichiarata
con Decreto del Prefetto di questa Congregazione per i Vescovi in data
1\textsuperscript{o} luglio 1988''). Bounded negative: the decree itself
is absent from the 1988 \work{Acta} volume, whose only Lefebvre entries
are at pages 1121 and 1495 (index checked 2026-07-25).}
\item On \textbf{2 July 1988} John Paul II issued the motu proprio
\work{Ecclesia Dei}, which is not primarily a penal act at all. It
declares the act schismatic, records the penalty, diagnoses its doctrinal
root, and then institutes a commission and orders a generous application
of the 1984 indult.
\end{itemize}

\subsection{What \work{Ecclesia Dei} decided, in its own words}

The motu proprio's third paragraph states the offence and the penalty
with precision: in itself the act ``was one of disobedience to the Roman
Pontiff in a very grave matter and of supreme importance for the unity of
the church, such as is the ordination of bishops whereby the apostolic
succession is sacramentally perpetuated. Hence such disobedience---which
implies in practice the rejection of the Roman primacy---constitutes a
schismatic act.'' The footnote to that sentence is canon 751, the
definition of schism; the footnote to the excommunication is canon 1382
of the 1983 Code as it then stood, which punished a bishop consecrating
without a pontifical mandate, and the one who receives consecration from
him, with excommunication reserved to the Apostolic See.\endnote{\work{Ecclesia
Dei}, n.~3 with footnotes 3 and 4 (\work{Acta Apostolicae Sedis} 80
(1988) 1496): ``Cf.\ Codex Iuris Canonici, can.~751''; ``Cf.\ ibid.,
can.~1382.'' The English quoted in the text is the Holy See's own English
version. The 1983 Code's canon 1382 was renumbered 1387 in the Book VI
revised by \work{Pascite gregem Dei} (2021); the substance of the offence
is unchanged, and the 2026 decree cites the new number. This account
never projects the new numbering backward.}

Paragraph 4 gives the doctrinal diagnosis: the root of the schismatic act
lies in ``an incomplete and contradictory notion of Tradition''---incomplete
because it does not sufficiently account for the living character of
Tradition as \work{Dei Verbum} 8 describes it, contradictory because it
sets Tradition against the universal Magisterium of the Roman Pontiff and
the body of bishops.

Paragraph 5~c) states the consequence for others, and it is the sentence
that has governed the status of the Society's adherents ever since:
``Everyone should be aware that formal adherence to the schism is a grave
offence against God and carries the penalty of excommunication decreed by
the Church's law,'' with a footnote to canon 1364---the canon on apostasy,
heresy, and schism.

Paragraph 6 institutes the commission, composed of a cardinal president
and other members of the Roman Curia, and orders that respect be shown
everywhere for those attached to the Latin liturgical tradition ``by a
wide and generous application of the directives already issued some time
ago by the Apostolic See for the use of the Roman Missal according to the
typical edition of 1962''---the 1984 indult again.

\actrecord{\work{Ecclesia Dei}, 2 July 1988 (\work{Acta Apostolicae
Sedis} 80 (1988) 1495--1498; Holy See English text).}{That the
consecrations of 30 June 1988 were performed after a formal canonical
warning of 17 June; that Lefebvre and the four named priests thereby
incurred excommunication under the then canon 1382; that the act is
declared schismatic under canon 751; that formal adherence to the schism
carries excommunication under canon 1364; and that a Roman commission was
instituted with a mandate of reconciliation and liturgical generosity.}{A
papal act declaring and diagnosing. It does not adjudicate the 1975
suppression, does not define ``formal adherence''---the 1996 note was
written precisely because bishops asked what it meant---and imposes no
penalty on the Society as an institution, which as a body without
canonical status could not be the subject of one.}

\subsection{The distinction that will matter later}

Read together, the acts of 1 and 2 July 1988 make three distinct
statements that later commentary constantly merges:

\begin{enumerate}
\item Six named individuals incurred a \latin{latae sententiae}
excommunication, declared on 1 July.
\item The act itself was characterized as schismatic.
\item Persons who \emph{formally adhere} to the schism incur
excommunication by the general law of canon 1364---a statement of law
about a class, not a declaration about any individual.
\end{enumerate}

In 2009 the first of these was undone for four men, and neither the
second nor the third was touched. In 2026 all three were repeated in a
harder form. Nothing in the intervening decades altered the framework laid
down in these two July days.
